\section{Metric-cone shortcuts for thin trees}\label{tg:section}

We next develop the geometric input to the thin-tree algorithm.  The
shortcut theorem of the companion \cite{ThinTreesCompanion} tests a matrix on cuts.  Convex
optimization in the next section will use a larger, finitely described
cone of tests.  We establish the corresponding hierarchy estimate by
extending its complete geometric arguments, which follow the hierarchy
and effective-resistance framework of Anari and Oveis Gharan
\cite[Sections~5--7]{AnariOveisGharan2015}.

Throughout this section graphs are finite loopless multigraphs on a
vertex set $V$, with $n=|V|\ge2$; edge sets retain multiplicities.
All symmetric matrices act on the Euclidean space
$\mathsf E=\one^\perp\subset\R^V$.  Equivalently, they are centered full
matrices, with their zero action on $\Span\{\one\}$ understood.
Write $\Id_{\mathsf E}$ for the identity on this space,
$b_{uv}=e_u-e_v$, and $L_G=\sum_{e\in E(G)}b_eb_e^T$, using arbitrary
edge orientations.  Matrix pairings are $\ip{A}{Z}=\tr(AZ)$.
Define the closed convex cone
\begin{equation}\label{tg:test-cone}
 \mathcal N=\left\{Z\succeq0:
 b_{uw}^TZb_{uw}\le b_{uv}^TZb_{uv}+b_{vw}^TZb_{vw}
 \quad(u,v,w\in V)\right\}.
\end{equation}
Its trace-dual cone is
$\mathcal N^*=\{A:\ip{A}{Z}\ge0\text{ for all }Z\in\mathcal N\}$.
The Gram-matrix representation of squared Euclidean distances is
classical \cite{Schoenberg1935}; the additional triangle inequalities
give the standard $\ell_2^2$ metric relaxation
\cite[Section~2]{AroraRaoVazirani2009}.
For $Z\in\mathcal N$, the numbers $d(u,v)=b_{uv}^TZb_{uv}$ form
a squared Euclidean pseudometric.  The centered identity belongs to
$\mathcal N$.  So does $c_Sc_S^T$, where
$c_S=\one_S-(|S|/n)\one$: its squared distances are precisely the
$0$--$1$ cut metric of $S$.

A \emph{hierarchy} is a rooted tree of distinct nonempty subsets of
$V$, with root $V$, singleton leaves, and children partitioning their
parent into at least two sets.  For a nonroot node $A$, with parent
$A^*$, set
\[
 P(A)=\delta_G(A),\qquad O(A)=E_G(A,A^*\setminus A).
\]
The hierarchy is \emph{locally $h$-connected} if every $G[A]$ is
$h$-edge-connected, with the singleton condition vacuous.
Each $|O(A)|\ge h$.  Each graph edge belongs to exactly two parent
boundaries: those of the two children of its endpoints' least common
ancestor.  In particular its multiplicity in any selected collection
of parent boundaries is at most two.

\Needspace{4\baselineskip}
\begin{theorem}[Hierarchy bound for the metric test cone]
\label{thm:tg-hierarchy}
There is an absolute integer $h_0$ such that, for every integer
$h\ge h_0$, every locally $h$-connected hierarchy of $G$, and every
collection $\mathcal M$ of nonroot nodes satisfying
$|O(A)|\ge |P(A)|/h$, there is a matrix $D\succ0$ on $\mathsf E$ with
\begin{align}
 \ip{D}{Z}&\le\ip{L_G}{Z} &&(Z\in\mathcal N),\label{tg:shortcut}\\
 \frac1{|O(A)|}\sum_{e\in O(A)}b_e^TD^{-1}b_e
 &\le h^{-31/40} &&(A\in\mathcal M).\label{tg:hierarchy-resistance}
\end{align}
The cutoff is independent of the graph and the hierarchy depth.
\end{theorem}

The proof first expresses the shortcut objective as a ratio whose
denominator is total edge length. We then represent its numerator by
collections of disjoint balls and charge their radii to edge lengths;
the resulting estimate will give simultaneous responses to partition
chains.

\subsection{The conic dual}

\begin{lemma}[A geometric upper bound for the shortcut program]
\label{tg:dual}
Let $G$ be connected and $\mathcal M$ a nonempty finite collection
of nonroot nodes with nonempty parent boundaries.  Let $p_*$ be the infimum, over
$D\succ0$ satisfying \Cref{tg:shortcut}, of the largest marked
average resistance.  Then
\begin{equation}\label{tg:dual-ratio}
 p_*\le \sup_{p,U}
 \frac{\displaystyle\sum_{A\in\mathcal M}\frac1{|O(A)|}
       \left(\sum_{e=uv\in O(A)}a_e\right)^2}
      {\displaystyle\sum_{e=uv\in E(G)}d(u,v)},
 \qquad
 a_e=|\ip{u_e}{p_u-p_v}|.
\end{equation}
Here $p_v$ are Euclidean points whose squared distances
$d(u,v)=\norm{p_u-p_v}_2^2$ satisfy the triangle inequalities, the
denominator is positive, and $U$ has orthonormal rows $u_e^T$.
The Euclidean dimension may be enlarged.
\end{lemma}

\begin{proof}
Put $A_A=|O(A)|^{-1}\sum_{e\in O(A)}b_eb_e^T$ and
$f_A(D)=\tr(A_AD^{-1})$.  These functions are convex, since
\[
 b^TD^{-1}b=\sup_x(2b^Tx-x^TDx).
\]
The conic constraint is $L_G-D\in\mathcal N^*$.
Choose $D_0=\delta\Id_{\mathsf E}$ with
$0<\delta<\lambda_{\min}(L_G|_{\mathsf E})$.
Then $L_G-D_0\succ0$ lies in the interior of $\mathcal N^*$,
because the positive semidefinite cone is contained in
$\mathcal N^*$.  Taking an epigraph height above all $f_A(D_0)$
makes the other inequalities strict as well.

For completeness, finite-dimensional separation gives the needed
strong duality even if the infimum is not attained.  Separate
$(0,0,p_*)$ from the open convex set of $(a,M,z)$ for which some
$D\succ0$ and $t\in\R$ satisfy
\[
 f_A(D)-t<a_A,\qquad
 M-D+L_G\in\operatorname{int}\mathcal N^*,\qquad t<z.
\]
Monotonicity in the three coordinates makes the separating
coefficients $(\lambda,Z,\gamma)$ satisfy
$\lambda\ge0$, $Z\in\mathcal N$, and $\gamma\ge0$.
The strictly feasible point just constructed rules out $\gamma=0$:
one may take all $a_A<0$ and $M$ in $-\operatorname{int}\mathcal N^*$,
making the pairing strictly negative for every nonzero
$(\lambda,Z)$.  Normalize $\gamma=1$.  The unrestricted variable $t$
then forces $\sum_A\lambda_A=1$.  Taking limits to the boundary of
the open set, and using weak duality for the reverse inequality,
gives
\begin{equation}\label{tg:lagrange}
 p_*=\sup_{\substack{\lambda\ge0,\ \sum_A\lambda_A=1\\Z\in\mathcal N}}
 \left[\inf_{D\succ0}
 \{\tr(A(\lambda)D^{-1})+\tr(ZD)\}-\tr(ZL_G)\right],
 \quad A(\lambda)=\sum_A\lambda_AA_A.
\end{equation}

For any $A,Z\succeq0$,
\begin{equation}\label{tg:matrix-infimum}
 \inf_{D\succ0}\{\tr(AD^{-1})+\tr(ZD)\}
 =2\tr\bigl((Z^{1/2}AZ^{1/2})^{1/2}\bigr).
\end{equation}
When $A,Z\succ0$, set $H=Z^{1/2}DZ^{1/2}$ and
$C=Z^{1/2}AZ^{1/2}$.  The objective minus $2\tr(C^{1/2})$ is
$\norm{H^{-1/2}C^{1/2}-H^{1/2}}_{\mathrm F}^2$, and vanishes at
$H=C^{1/2}$.  For singular $A$ or $Z$, add
$\eta\Id_{\mathsf E}$ to both.  The positive definite identity at
each fixed $D$, followed by $\eta\downarrow0$, proves the lower
bound.  The unperturbed infimum is at most the perturbed infimum,
whose limit gives the upper bound.  These are auxiliary perturbations
for the identity, not extra dual feasibility assumptions.

For $Z\ne0$ in $\mathcal N$, connectivity gives
$\mathcal E=\tr(L_GZ)>0$.  Writing
$N_0=\tr((Z^{1/2}A(\lambda)Z^{1/2})^{1/2})$, optimization over the
scale $\tau Z$ gives
\[
 \sup_{\tau\ge0}(2\sqrt\tau N_0-\tau\mathcal E)
 =N_0^2/\mathcal E.
\]
The zero multiplier contributes zero.  Represent $Z$ as the centered
Gram matrix of points $p_v$, and put
\[
 w_e=\sum_{A:e\in O(A)}\frac{\lambda_A}{|O(A)|}.
\]
Then $N_0$ is the nuclear norm of the matrix with columns
$\sqrt{w_e}(p_u-p_v)$.  The singular value decomposition, after
adding zero coordinates if needed, gives
\[
 N_0=\max_{UU^T=I}\sum_e\sqrt{w_e}\ip{u_e}{p_u-p_v}.
\]
Changing individual row signs preserves row orthonormality, so take
the displayed contributions nonnegative.  For such a maximizing $U$,
\begin{align*}
 N_0
 &\le\sum_A\sqrt{\frac{\lambda_A}{|O(A)|}}
                  \sum_{e\in O(A)}a_e,\\
 N_0^2
 &\le\sum_A\frac1{|O(A)|}
                   \left(\sum_{e\in O(A)}a_e\right)^2.
\end{align*}
The first line uses $\sqrt{\sum_jx_j}\le\sum_j\sqrt{x_j}$ and
the second uses $\sum_A\lambda_A=1$.  Overlapping boundaries are
allowed.  Substitution into \Cref{tg:lagrange} proves the claim.
\end{proof}

\subsection{A convex realization with exact radial motion}

Starting from the classical distance-profile isometry into
$\ell_\infty$ \cite{Frechet1910,Kuratowski1935}, we construct a
continuous geometry for an arbitrary finite metric.  It will be used for balls and shells; the Euclidean
points $p_v$ remain the inputs to the spectral argument.
\Cref{tg:geometry-figure} contrasts the Euclidean spectral points with
the distance-profile realization of the same vertex labels and
illustrates the exact radial motion of \Cref{tg:profiles}.

\begin{lemma}[Doubled distance profiles]\label{tg:profiles}
Every finite metric $(S,d)$ has an isometric realization
$x\mapsto c_x$ in a convex polytope $Q$ in a finite-dimensional
normed space with the following properties.  Distance
$\rho_x(q)=\norm{q-c_x}$ is affine on $Q$.  There is a common $R>0$
and a point $c_x^-\in Q$ such that
\begin{equation}\label{tg:opposite-radii}
 \norm{q-c_x}+\norm{q-c_x^-}=R\qquad(q\in Q).
\end{equation}
From any $q\in Q$ one can reach every radial value in $[0,R]$ by
moving toward $c_x$ or $c_x^-$, and the change in radius equals
the distance moved.  Every straight segment in $Q$ has its norm
length, and clamped radial functions are piecewise affine on it.
\end{lemma}

\begin{proof}
Take $R>2\max_{x,y}d(x,y)$ and two copies $x^+,x^-$ of each point.
On the doubled set define
\[
 \widehat d(x^+,y^+)=\widehat d(x^-,y^-)=d(x,y),\qquad
 \widehat d(x^+,y^-)=R-d(x,y).
\]
This is a metric.  A triangle with all signs equal is an original
triangle.  A triangle with two equal signs has two cross-copy sides;
an inequality with a cross-copy side on the left reduces to an
original triangle inequality.  The remaining inequality follows
from $d(x,y)+d(x,z)+d(z,y)\le3\operatorname{diam}(d)<2R$.

For a doubled point $a$, let
$c_a=(\widehat d(a,z))_z$ be its distance profile in $\ell_\infty$,
and set $Q=\conv\{c_a\}$.  The triangle inequality and coordinate
$a$ show $\norm{c_a-c_b}_\infty=\widehat d(a,b)$.
More generally, for $q=\sum_b\theta_bc_b$,
\begin{equation}\label{tg:affine-radius}
 \norm{q-c_a}_\infty
 =\sum_b\theta_b\widehat d(a,b)=q_a,
 \qquad \theta_b\ge0,\quad\sum_b\theta_b=1.
\end{equation}
Indeed, convexity of the norm gives the upper bound, and coordinate
$a$ gives the lower bound.  Thus the radius is the restriction of
a linear coordinate function, independently of the representation
of $q$.  Set $c_x=c_{x^+}$ and $c_x^-=c_{x^-}$.
The identity
$\widehat d(x^+,b)+\widehat d(x^-,b)=R$ for every doubled point $b$
proves \Cref{tg:opposite-radii}.

If $\rho=\rho_x(q)$, the segment $q_t=(1-t)q+tc_x$ has
$\norm{q_t-q}=t\rho$ and $\rho_x(q_t)=(1-t)\rho$.
The segment $q_t^-=(1-t)q+tc_x^-$ has
$\norm{q_t^--q}=t(R-\rho)$ and
$\rho_x(q_t^-)=\rho+t(R-\rho)$.
These formulas prove the exact-motion assertion, including the
endpoint cases.  Convexity of $Q$ keeps all segments in $Q$;
the final assertion follows from \Cref{tg:affine-radius} and clamping.
\end{proof}

For a pseudometric, first identify zero-distance points to apply
\Cref{tg:profiles}.  This identification concerns only the geometric
point set.  Every graph vertex, edge copy, hierarchy node, and
center-vertex label is retained.  In particular coincident vertex
images do not contract the graph.

\begin{figure}[tb]
\centering
\begin{tikzpicture}[font=\small,>=Stealth]
 \node[font=\bfseries] at (2.7,2.6) {Two realizations of the vertex labels};
 \node[draw,rounded corners,align=center,minimum width=4.5cm,
       fill=MidnightBlue!5] (euclidean) at (2.7,1.55)
       {$p_v$ and $Y_v=Up_v$ in Euclidean space\\[3pt]
        $\norm{Y_u-Y_v}_2^2\le d(u,v)$};
 \node[draw,rounded corners,align=center,minimum width=4.5cm,
       fill=ForestGreen!5] (profiles) at (2.7,-.2)
       {$c_v$ in the profile polytope $Q$\\[3pt]
        $\norm{c_u-c_v}_\infty=d(u,v)$};
 \draw[->,thick] (euclidean) -- node[right,align=left,font=\footnotesize]
       {selected vertex labels} (profiles);
 \node[font=\bfseries] at (10.2,2.6) {Exact radial motion in $Q$};
 \coordinate (center) at (7.6,1.3);
 \coordinate (point) at (10.2,-.2);
 \coordinate (opposite) at (12.8,1.3);
 \draw[->,thick,MidnightBlue] (point) -- (center);
 \draw[->,thick,ForestGreen] (point) -- (opposite);
 \fill (center) circle (2pt);
 \fill (point) circle (2pt);
 \fill (opposite) circle (2pt);
 \node[anchor=south,align=center,font=\footnotesize]
       at (7.6,1.48) {$c_x$\\[2pt]radius $0$};
 \node[anchor=south,align=center,font=\footnotesize]
       at (12.8,1.48) {$c_x^-$\\[2pt]radius $R$};
 \node[anchor=north,align=center,font=\footnotesize,
       text=MidnightBlue] at (8.55,.03)
       {move inward\\[2pt]length $\rho$};
 \node[anchor=north,align=center,font=\footnotesize,
       text=ForestGreen] at (11.85,.03)
       {move outward\\[2pt]length $R-\rho$};
 \node[anchor=north,align=center,font=\footnotesize]
       at (10.2,-.4) {$q$\\[2pt]radius $\rho$};
\end{tikzpicture}
\caption{Spectral packing uses Euclidean points; balls and shells use
the distance-profile realization with the same vertex labels.
The radial paths are schematic: in the ambient norm, distance
traveled along either indicated segment equals change of radius.
The two segments need not form one straight segment in $Q$.}
\label{tg:geometry-figure}
\end{figure}


\subsection{Transferring the bag and charging estimates}

The next lemma extends the bag-selection and shell-charging arguments of
\cite[Propositions B.7, B.8, and B.10]{ThinTreesCompanion} to $Q$.
We retain those constructions and their scalar accounting, but verify
the geometric steps below: their stated cube hypotheses do not cover
all tests in $\mathcal N$.  The two charging alternatives explain the
division into compact and assigned bags.  Compact bags charge disjoint
geometric objects.  Assigned bags may overlap geometrically, provided
their hierarchy labels separate the graph edges used for charging.

Use relative open balls $B(c,\delta)=\{q\in Q:\norm{q-c}<\delta\}$
and strict shells $B(c;a,b)=\{q\in Q:a<\norm{q-c}<b\}$.
An ordinary ball includes its center, whereas a shell with $a=0$
excludes it.  Both are called radial objects.  An ordinary ball
$B(c,b)$ has \emph{numerical} inner radius $a=0$ and width $b$;
a shell has width $b-a$.  All radial tests and width formulas below
apply to both types with these numerical radii, but membership and
intersection always use the actual support.  In particular, an
unchanged ordinary ball keeps its center.  A radial restriction to
$a\le a'<b'\le b$ means replacement by the strict shell
$B(c;a',b')$, even when $a'=0$.  Such a replacement is a subset
of the old support; it is an operation, not an identification of a
ball with a punctured shell.  Proposed intervals of nonpositive
width are discarded.
A \emph{bag} consists of mutually disjoint balls of
a common radius $\delta$, centered at vertex images $c_v$.
It is \emph{$\beta$-compact of type $(\delta,\Delta)$} if it has
at least two balls, its center diameter is at most $\Delta$, and
$\beta\Delta\le\delta|\mathrm{Bag}|$.
It is \emph{$\eta$-assigned to $A$} if
$|\mathrm{Bag}|\ge\eta|O(A)|$ and every ball has a designated center
vertex $u\in A$ with an edge $uv\in O(A)$ of length $d(u,v)<\delta$.
A family has all its balls mutually disjoint and a common radius.
A compact family has a common type $(\delta,\Delta)$.
An assigned family has one bag for each of a set of distinct assigned
nodes. Bags assigned to the same node at one scale may always be
merged: this preserves every witnessing edge, the assignment bound,
and the total mass.  Fix $0<\rho<1$. A sequence of compact
families is $\rho$-geometric when
$\Delta_{i+1}<\rho\delta_i$; a sequence of assigned families is
$\rho$-geometric when $\delta_{i+1}<\rho\delta_i$ and its assigned
nodes are distinct across families.  In either case its mass is
\[
 \mathsf M=\sum_i\delta_i\sum_{\mathrm{Bag}\in\mathcal F_i}
                                      |\mathrm{Bag}|.
\]

\begin{lemma}[Geometric transfer]\label{tg:transfer}
Let $p_v\in\R^{d_{\mathrm E}}$ be points for which
$d(u,v)=\norm{p_u-p_v}_2^2$ is a pseudometric.  Realize $d$ in $Q$
as in \Cref{tg:profiles}.  Let $U:\R^{d_{\mathrm E}}\to\R^{E(G)}$ be a
contraction, change its row signs so
$a_e=\ip{u_e}{p_u-p_v}\ge0$, and set
\[
 d_e=d(u,v),\qquad
 N_A=\frac{(\sum_{e\in O(A)}a_e)^2}{|O(A)|},\qquad
 \mathcal E=\sum_{e\in E(G)}d_e.
\]
The following three assertions hold in $Q$.
\begin{enumerate}[label=(\roman*)]
\item\label{tg:transfer-extract}
Let $T$ be a nonempty set of nodes with $N_A>0$ and
$N_A\ge\alpha\sum_{e\in O(A)}d_e$, where $0<\alpha\le1$.
Define
\begin{align*}
 J&=1+\lceil\log_2(256/\alpha^2)\rceil,&
 C_2&=128J^4,&\sigma&=\alpha/C_2,\\
 L&=2+\lceil\log_2(1/\sigma)\rceil,&
 H_0&=1+\lceil\log_2(24/(\rho\sigma^2))\rceil,&
 C_1&=1024\epsilon^{-2}.
\end{align*}
If $\beta>1$, $0<\epsilon<1/3$, $0<\rho<1$, and
$\sigma^\epsilon\le(48\beta C_1)^{-1}$, there is a
$\rho$-geometric sequence consisting either of $\beta$-compact
families or of $\sigma^{1+2\epsilon}$-assigned families, with
\begin{equation}\label{tg:captured-mass}
 \mathsf M\ge
 \frac{\sigma^\epsilon}{384\beta C_1C_2LH_0}
                  \sum_{A\in T}N_A.
\end{equation}
\item\label{tg:transfer-compact}
For $G$ $k$-edge-connected, a $\rho$-geometric sequence of $\beta$-compact families
with $\rho\le1/12$ and $\beta\ge36$ satisfies
$\mathcal E\ge(k/4)\mathsf M$.
\item\label{tg:transfer-assigned}
For a locally $k$-connected hierarchy, selected nodes satisfying
$|O(A)|\ge k\rho|P(A)|$, and a $\rho$-geometric sequence of
$(24C_3/k)$-assigned families, one has
\begin{equation}\label{tg:assigned-charge}
 \mathcal E\ge\frac{kC_4}{96C_3}\mathsf M
\end{equation}
provided
$C_4\ge3$, $\rho\le(6C_4)^{-1}$, and
$C_3\ge2((C_4+1)+4(C_4+2)^2)$.
\end{enumerate}
\end{lemma}

\begin{proof}
We follow the constructions in Sections B.2--B.4 of
\cite{ThinTreesCompanion}, keeping their scale choices and scalar
potentials.  All balls, shells, and intersections now lie in the
continuous polytope $Q$.  New shells inherit vertex centers, so
\Cref{tg:affine-radius} remains available.  We first obtain the bags,
then explain how each of the two charging constructions operates in
this realization.

\emph{Extraction and the spectral step.}
The two truncations and homogeneity arguments in
\cite[Section B.2.2]{ThinTreesCompanion} use only
$a_e^2\le d_e$, nonnegative boundary sums, and multiplicity at most
two.  They select nonempty subsets $O'(A)\subseteq O(A)$ retaining
a fixed fraction of each squared boundary sum, and group the nodes
so that both positive quantities $a_e^2$ and $d_e$ vary by a factor
less than four within $F=\bigcup_AO'(A)$.  For one such group put
$N=\sum_A|O(A)|$, $N'=|F|$,
$c_1=\min_{e\in F}a_e^2$, and $c_2=\max_{e\in F}d_e$.
The truncation and grouping estimates give $c_1N'\ge\sigma Nc_2$.
For the Euclidean points $Y_v=Up_v$ and
$\widetilde\alpha=\sigma N/N'$ this implies
\[
 (\E_{e\in F}a_e)^2\ge c_1
 \ge\widetilde\alpha c_2
 \ge\widetilde\alpha\E_{e\in F}\norm{Y_u-Y_v}_2^2.
\]
This is the exact hypothesis of the Euclidean spectral-packing
Lemma B.3 of \cite{ThinTreesCompanion}; that lemma places no binary
restriction on its points.  Its selected centers have squared
Euclidean separation at least $4r$.  Since
\begin{equation}\label{tg:spectral-to-metric}
 \norm{Y_u-Y_v}_2^2\le d(u,v)=\norm{c_u-c_v}_\infty,
\end{equation}
the corresponding $Q$-balls of radius $r$ are disjoint.  If $b_0$
is their number, the spectral lemma gives
$b_0r\ge\widetilde\alpha^\epsilon c_1N'/C_1$.  Shrink their common
radius to $\delta$ defined by
$b_0\delta=\widetilde\alpha^\epsilon c_1N'/C_1$.
Only the selected vertex labels pass from the Euclidean geometry
to $Q$; a map on intermediate points is unnecessary.

The rest of the one-family construction uses a maximal disjoint
collection of balls of radius $D=\max\{\delta,2c_2\}$, centered at
images $c_u$ for which $u\in A$ and $uv\in O'(A)$ for some $A$.
Maximality puts every candidate center within $2D$ of a selected
center: intersection of two relative balls implies this distance
bound by the triangle inequality.  Every endpoint center is within
$c_2$ of an assigned endpoint, hence within $3D$ of a selected center.
Grouping by those centers gives diameter at most $6D$ in the compact
case.  In the assigned case the witnessing edge has length
at most $c_2<D$.  The cardinality estimates of
\cite[Lemma B.6]{ThinTreesCompanion} give compactness, assignment
strength, and radius range
$[\sigma c_1,12c_1/\sigma]$.  Its final grouping by dyadic size and
then by residue class modulo $H_0$ is numerical: each scale lies
between $\sigma2^j$ and $24\,2^j/\sigma$, so scales in the same
residue class meet the stated geometric separation.  The losses
$C_2$, $L$, $2H_0$, and $192\beta C_1/\sigma^\epsilon$ are exactly
those in \Cref{tg:captured-mass}.  This proves (i).

\emph{Length charging and compact radial objects.}
The resource charged by both arguments is one fixed straight
segment $[c_u,c_v]$ of length $d_e$ for each graph edge, retaining
parallel edges as separate resources.  For a radial object with numerical radii
$a<b$, its clamped radial function is
$f(q)=\min\{b,\max\{a,\norm{q-c}\}\}$.
It is $1$-Lipschitz and piecewise affine on each edge segment.
If the object has vertex witnesses at radius at most $a$ and at
least $b$, all its intermediate radial cuts are nontrivial.
Consequently
\[
 k(b-a)\le\int_a^b|\delta_G\{v:\norm{c_v-c}<t\}|\,dt
 =\sum_{uv\in E(G)}|f(c_u)-f(c_v)|.
\]
For disjoint radial objects, subdivide an edge segment at the finitely many
breakpoints of these functions.  At most one clamped function is
nonconstant on each open subsegment, and its variation is bounded
by that subsegment's length.  The total charge is therefore at most
$d_e$, including edges along a boundary or of zero length.  This
proves the length bound of \cite[Lemma B.9]{ThinTreesCompanion} in $Q$.
For an ordinary ball, nonconstant portions of the clamped function
lie at strictly positive radius.  Its center contributes no charged
length, but remains part of the ball in membership tests.

For compact bags, process the families in decreasing order of radius,
as in \cite[Proposition B.8]{ThinTreesCompanion}.  An input ball inserted
at the current scale is called fresh and is left intact until the
next scale; all other live objects are called old.  Thus a fresh
ordinary ball becomes an old ordinary ball at the next scale unless
a subsequent operation restricts it.  Disjoint
equal-radius balls have centers separated by at least $2\delta$,
since their midpoint belongs to $Q$.  For any live object with center
$c$ and numerical radii $a,b$, call $B(c_u,\delta)$ interior when
\[
 a+\delta+\Delta\le\norm{c_u-c}\le b-\delta-\Delta.
\]
This conservative two-sided test also applies when the object is
an ordinary ball, with $a=0$.  If one ball of a compact bag is
interior, its containing object is old, since fresh balls and
unprocessed balls belong to the same disjoint family.  All bag centers
have radial values in an interval $[l,u]\subseteq[a+\delta,b-\delta]$
of width at most $\Delta$.  Replacing the old object by the strict
shells with intervals $(a,l-\delta)$ and $(u+\delta,b)$ leaves room
for all the new balls, including when the old object was ordinary.
The extremal centers supply the new outer and inner witnesses;
the previous witnesses supply the other two.  Every inserted ball
has its own center and a different bag center as witnesses.
All supports are disjoint by $1$-Lipschitzness of radius.
When no interior bag remains, shrinking each old interval at both
ends by $2\delta+\Delta$ separates it from every remaining ball.
Indeed, failure at the lower margin puts the entire candidate at
radius less than $a+2\delta+\Delta$, and failure at the upper
margin puts it beyond $b-2\delta-\Delta$.  This includes candidates
near the center of an old ordinary ball: the shrink removes that
center neighborhood.  Witnesses are preserved because the inner
radius increases and the outer radius decreases; an inner witness
need not belong to the remaining strict shell.
The capacities $[w-6\Delta_i]_+$ for all old objects and
$\delta_i-6\Delta_{i+1}$ for fresh balls depend only on widths and
the inequalities $2\delta_i\le\Delta_i$,
$6\Delta_{i+1}\le\delta_i/2$, and
$|\mathrm{Bag}|\delta_i\ge36\Delta_i$, all verified here.
The split loses numerical width $(u-l)+2\delta$, and the shrink
loses $2(2\delta+\Delta)$, with the same positive-part accounting
when a residual is discarded.  Removing the center alone adds no
width loss.  A retained fresh ball has the same capacity after
becoming old.  Thus the split and scale-transition inequalities in
the cited proof retain at least one quarter of the input mass as
final radial-object width.
Applying the preceding length bound proves (ii).

\emph{Assigned labels, pruning, and path charging.}
Assigned bags require the labeled records of
\cite[Section B.4.1]{ThinTreesCompanion}.  We give their definitions
because geometric overlap is allowed.  Write $\mathcal T[A]$ for
the hierarchy subtree rooted at $A$, identifying its leaves with the
vertices in $A$.  A record $H$ has an ordinary ball or shell as its
support, an owner $A_H$, a set $\mathcal P_H$ of strict descendants
of $A_H$, and optionally one avoided strict descendant $D_H$.
Define its conflict domain and leaf set by
\[
 C(H)=\mathcal T[A_H]\setminus
 \left(\bigcup_{P\in\mathcal P_H}\mathcal T[P]
       \ \cup\ \mathcal T[D_H]\right),
 \qquad L(H)=C(H)\cap V,
\]
omitting the final subtree when there is no avoided label.
The record is called avoiding precisely when $D_H$ is present.
For center $c$ and numerical radii $a<b$, validity means:
\begin{enumerate}[label=(\alph*)]
\item the labels have the stated form, so $A_H\in C(H)$;
\item both endpoints of every edge of $\delta_G(P)$,
      for every $P\in\mathcal P_H$, have radius at least $b$ from $c$;
\item some $u\in L(H)$ has radius at most $a$ and starts at least
      $k/4$ edge-disjoint paths to vertices of radius at least $b$.
      These paths lie in $G[A_H]$, or in $G[A_H\setminus D_H]$
      for an avoiding record;
\item geometrically intersecting records have disjoint conflict domains.
\end{enumerate}
The initial paths in (c) need not avoid the predecessor subtrees
in $\mathcal P_H$; condition (b) will remove those subtrees from
the portions used for charging.  Its inner witness is a graph
vertex in $L(H)$, not necessarily a point of the support: for a
strict shell it can lie on or inside the excluded inner boundary.

For a node $A$ assigned at scale $\delta_i$, its predecessors are
the strict descendants assigned at earlier, larger scales.  An input
ball is noninsertable if it contains an endpoint of a boundary edge
of any predecessor.  Otherwise give it owner $A$, no avoided label,
and excluded set equal to the predecessors.

At one scale, disjoint center-inclusive balls contain at most
$2|P(A)|$ endpoint vertices of $P(A)$.  Thus a predecessor $A$,
whose own scale is $\delta_i$, deletes mass at most
$4\rho\delta_i|P(A)|$, which the boundary ratio and
assignment bound pay as in the source's Lemma B.13.  Deleting
noninsertable balls and bags losing more than half their balls
retains mass at least $\mathsf M/2$ and assignment at least
$12C_3/k$.  Insertability puts the designated center vertex in
its conflict domain: otherwise its witnessing edge in $O(A)$
would be a forbidden predecessor-boundary edge whose endpoint
lies in the ball.  Another ball of the bag supplies an outer
vertex; local connectivity supplies the required paths.
The retained bag has more than one ball, since
$|O(A)|\ge k$ and its assignment is at least $12C_3/k$.
Thus the resulting input records satisfy (a)--(c), including when
distinct vertices have coincident images.

For a valid collection, stop each witnessing graph path at its
first outer vertex.  It cannot enter an excluded predecessor
subtree earlier, because the entering edge would have an endpoint
inside the outer radius.  On the resulting continuous path keep
the portion from the last inner-radius point to the first
outer-radius point.  Its length is at least the record's numerical
width by $1$-Lipschitzness.  Apart from endpoints it lies in
$a<\norm{q-c}<b$, which is contained in either type of support.
For an ordinary ball $a=0$ and the center carries no charged length.
Paths within one record are edge-disjoint.
Both endpoints of every edge in a stopped path belong to $L(H)$:
an entrance into an excluded subtree, even at the stopping vertex,
would have a preceding endpoint inside the outer radius, contradicting
(b).  Two records charging the same graph edge therefore have intersecting
conflict domains, so (d) makes their actual supports disjoint and
they charge disjoint open portions of that edge.  This proves
the source's Lemma B.14 in $Q$, at rate $k/4$.

It remains to construct a valid arrangement with enough total width.
Start with the empty arrangement and process scales in decreasing
order and, within each scale, bags in
increasing depth of their assigned nodes.  As in the compact
construction, original balls inserted during the current phase are
fresh and remain intact until the next phase; all other records are
old.  A candidate
$J=B(c_u,\delta_i)$ is interior to a current record
$H$ with center $c$ and numerical radii $a,b$ when
\[
 C(J)\cap C(H)\ne\varnothing,
 \qquad a+C_3\delta_i<\norm{c_u-c}<b-C_3\delta_i.
\]
A candidate interior to no current record is called a border
candidate.  In particular old ordinary balls participate in this
test with $a=0$, including its lower margin.  A record containing an
interior candidate has width greater than $2C_3\delta_i$, so it
cannot be a fresh ball of width $\delta_i$.  An interior candidate is
contained in $H$.  The processing order implies that the owner of
any conflicting current record is an ancestor of the candidate's
owner $A$, and that $A\in C(H)$.  Two conflicting records therefore
cannot both contain this candidate, by (d).  The strict concentric
restrictions defined above retain (a)--(c) with the old labels:
the old inner witness still has radius at most the new inner
radius, and the old paths still reach the smaller outer radius.
This remains true when an ordinary ball is replaced by a shell
and its center witness no longer lies in the support.
Restricting both support and conflict domain
preserves (d).  Discard any proposed radial interval of nonpositive
width.  Such restrictions cannot turn a border candidate
into an interior candidate.  These are the order and inheritance
properties of \cite[Lemmas B.17--B.20]{ThinTreesCompanion}.

\emph{Meeting a record and carving its interior.}
The following consequence of \Cref{tg:profiles} identifies the record
affected by a carving operation.  If a live record $H$ with center
$c$ and numerical radii $r_1<r_2$, and a vertex center $z$, satisfy
\[
 r_1-\delta<\norm{z-c}<r_2+\delta,
\]
then $B(z,\delta)$ meets the actual support of $H$.
Choose a radius in $(r_1,r_2)$
less than $\delta$ away from $\norm{z-c}$ and use the exact inward
or outward motion to reach it.  Its traveled distance is less than
$\delta$.  There is no radius-range obstruction: condition (c)
provides an outer vertex witness, so $r_2\le\operatorname{diam}(d)<R$.
For an ordinary ball $r_1=0$; the chosen radius is positive and
belongs to that ball, even when $z=c$ and outward motion is needed.
This replaces motion toward the binary center or its opposite cube
vertex in \cite[Section B.4.3]{ThinTreesCompanion}.

At scale $\delta$, process the interior candidates of a bag assigned
to $A$ while at least half its original balls remain interior and
untreated.  Candidates that cease to be interior are kept as border
candidates.  Carving an avoiding shell around an interior candidate
replaces its radial interval $(r_1,r_2)$ by
$(r_1,q-\delta)$ and $(q+\delta,r_2)$, and inserts the candidate.
Here $q$ is its center's radius from the shell center.
Every avoiding record is indeed a strict shell: input records are
nonavoiding, only the central shell $H_4$ below introduces an
avoided label, and later restrictions are strict shells.
Radial Lipschitzness gives disjointness; witnesses and labels
survive in the restrictions.  If no remaining interior candidate
belongs to an avoiding shell, the other two operations first split
$G[A^*\setminus A]$ into $j\le2|O(A)|/k$ induced subgraphs
that are $k/4$-edge-connected.  This count follows by successive
cuts smaller than $k/4$ and the $k$-connectivity of $G[A^*]$;
it is independent of the embedding.  Assignment and the
half-bag loop condition give one piece adjacent along edges of
length less than $\delta$ to at least $3C_3$ candidate centers.
Radial coordinates of these neighbors differ by less than
$\delta$ from those of their candidate centers.

Choose a nonavoiding record $H$ containing one of those candidates
in its interior, with center $c$ and numerical radii $r_1,r_2$,
and write $\varrho(v)=\norm{c_v-c}$.  It may be an old ordinary
ball, in which case $r_1=0$.
Let $[a,b]\subseteq[r_1,r_2]$ be the radial range of the neighbors
in the selected piece, clipped to this numerical interval.  Select for treatment
the $m$ still-interior candidates whose centers have radii in
$(a-\delta,b+\delta)$.  Each meets the chosen record by the preceding
exact-motion argument.  All candidates of the bag have the same
labels, so each has conflict domain meeting $C(H)$.  If such a
candidate is interior to another record $H'$, its entire ball lies
in $H'$, so $H$ and $H'$ intersect.  The processing-order property
puts its owner $A$ in both conflict domains, contrary to (d).
Thus every selected candidate is interior to $H$ itself.
The radial count in the cited proof gives
$m\ge3(C_3-1)$ in the dense case below and
$b-a\ge(C_3-1)\delta$ in the sparse case.  It uses only the
$3C_3$ adjacent centers, their interior margins, and the fact that
neighbor radii differ by less than $\delta$.

For the dense case, with $m\delta>3(b-a)$, retain intervals
$(r_1,a-2\delta)$ and $(b+2\delta,r_2)$ and insert the $m$ balls.
The expanded radial interval and Lipschitzness give every required
separation.  For the sparse case $m\delta\le3(b-a)$, retain
$(r_1,a)$ and $(b,r_2)$ and create the central shells
$H_3=B(c;a+\delta,b-\delta)$ and $H_4=B(c;a,b)$.
Give $H_3$ owner $A$, no avoided label, and the old excluded nodes
strictly below $A$.  Give $H_4$ the old owner and excluded set,
and avoided label $A$.  Their domains are respectively
$C(H)\cap\mathcal T[A]$ and $C(H)\setminus\mathcal T[A]$.
The old owner is a strict ancestor of $A$.  An original ball with
owner $A$ cannot come from an earlier scale, since assigned nodes
are distinct across families; at the current scale it is disjoint
from every unprocessed candidate.  A synthetic shell with owner
$A$ can only have been created while processing this same bag.
It has already treated every then-interior candidate in its expanded
interval; later restrictions only reduce its support, and border
candidates never become interior.  Hence neither an original ball
nor a synthetic shell with owner $A$ can be the chosen record.
Thus both new labels have the required form.
For the changed-label path witnesses, extremal neighbors
$v_1,v_2$ in the selected piece and their adjacent candidates
$u_1,u_2\in A$ satisfy
\[
 \varrho(v_1)\le a,\quad \varrho(v_2)\ge b,\qquad
 \varrho(u_1)<a+\delta,\quad \varrho(u_2)>b-\delta.
\]
Connectivity of $G[A]$ supplies $k$ paths for the inner central
shell; connectivity of the selected piece supplies at least
$k/4$ paths for the avoiding shell.  Their initial vertices
avoid excluded subtrees: membership in one would make the
short edge $u_1v_1$ an excluded boundary edge with an endpoint
inside the old outer radius, or make that excluded subtree an
ancestor of $A$, contradicting $A$'s conflict membership.
The sparse-case width bound gives $a+\delta<b\le r_2$, so the
initial vertices used here lie strictly inside the old outer radius.
This verifies (c) for both new labels; (b) follows from the smaller
outer radii and excluded sets.  Other records inherit separation by support and
conflict-domain inclusion.  Remaining interior candidates have centers
outside $(a-\delta,b+\delta)$, so miss both central shells.
If $a=0$, $H_4$ excludes its center but may still use $v_1$ at
radius zero as its inner witness, as condition (c) permits.

\emph{The scalar accounting and the phase transition.}
The potentials in Section B.4.3 are
$\delta_i-C_4\delta_{i+1}$ for fresh balls,
$[w-C_4\delta_i]_+$ for other nonavoiding records, and the
latter divided by $2(C_4+2)$ for avoiding records.
An original ball becomes old when its phase ends and then uses the
second formula with width equal to its radius; if retained unchanged,
it still has its center-inclusive support.  These functions use the
same numerical widths and radii in $Q$.  The three
replacements have formal width losses $2\delta$ across the two avoiding residual
shells in the first operation, $(b-a)+4\delta$ across the two
nonavoiding residual shells in the second, and $2\delta$ across
the three nonavoiding pieces in the third.  Its numerical
inequalities preserve old credit and supply
$C_4\delta/(6C_3)$ per treated candidate, counting sparse-case
candidates as treated even though their balls are not inserted.
Finally, Section B.4.4 either retains every record when the gain
from changing scale supplies the required border credit, or shrinks
every old object, including old ordinary balls, at both ends by
$(C_3+1)\delta_i$ and inserts remaining border balls when the
gain is insufficient.  For a conflicting old object, border status
means the candidate center has radius at most $a+C_3\delta_i$
or at least $b-C_3\delta_i$.  Its whole ball therefore lies below
the new inner radius or above the new outer radius, respectively.
This proves separation also at the center of an old ordinary ball;
fresh balls are disjoint within the current family.  Restrictions
retain witnesses.  Each old width decreases by
$2(C_3+1)\delta_i$, unless the object is discarded, while the
potential threshold decreases from $C_4\delta_i$ to
$C_4\delta_{i+1}$.  These are the same changes for a ball and a
shell.  The source's width
calculation bounds the lost potential by half the inserted
mass, leaving at least $C_4/(6C_3)$ of each border ball's radius
for its node.  The half-bag stopping rule leaves final width at
least $C_4/(12C_3)$ times the retained mass.  Combined with the
pruning factor $1/2$ and the path-charging rate $k/4$, this is
exactly \Cref{tg:assigned-charge}.

Thus the inherited constructions give (i)--(iii), with the same
scalar losses and with every length charge measured by
$\mathcal E=\sum_e d_e$ in the profile realization.
\end{proof}

\subsection{Choosing parameters and responding to a partition chain}

\begin{proof}[Proof of \Cref{thm:tg-hierarchy}]
If $\mathcal M$ is empty, a sufficiently small positive centered
identity satisfies \Cref{tg:shortcut}.  Otherwise use
\Cref{tg:dual}.  Fix one of its embeddings and row-orthonormal
systems, with $\mathcal E>0$, and set
\[
 \alpha=h^{-4/5},\qquad \epsilon=1/10,\qquad
 \beta=36,\qquad \rho=(2h^2)^{-1},\qquad
 C_3=10^4,\qquad C_4=3.
\]
Let $T$ consist of the marked nodes with positive $N_A$ and
$N_A\ge\alpha\sum_{e\in O(A)}d_e$.
Boundary multiplicity at most two gives
\begin{equation}\label{tg:good-nodes}
 \sum_{A\in\mathcal M\setminus T}N_A\le2\alpha\mathcal E.
\end{equation}
If $T$ is nonempty, apply \Cref{tg:transfer}(i).
Its parameters obey
\[
 C_2=O((1+\log h)^4),\quad L,H_0=O(1+\log h),\quad
 \sigma=h^{-4/5}/C_2.
\]
The smallness condition holds for all sufficiently large absolute
$h$.  So does the assignment requirement, because
\[
 h\sigma^{1+2\epsilon}
 =h^{1/25}/C_2^{6/5}\longrightarrow\infty,
 \qquad \sigma^{1+2\epsilon}\ge24C_3/h.
\]
The marked-node condition implies
$|O(A)|\ge|P(A)|/h\ge h\rho|P(A)|$.
For large $h$, $\rho\le1/18$; also
$10^4\ge2(4+4\cdot25)=208$.
Both charging conclusions of \Cref{tg:transfer} therefore apply,
according to the family type, and give the common bound
\[
 \mathcal E\ge c_*h\mathsf M,
 \qquad c_*=\frac{C_4}{96C_3}=\frac{1}{320000}.
\]
Together with \Cref{tg:captured-mass}, this yields
\[
 \frac{\sum_{A\in T}N_A}{\mathcal E}
 \le \frac{384\beta C_1C_2LH_0}{c_*h\sigma^\epsilon}
 =O\bigl(h^{-23/25}(1+\log h)^7\bigr).
\]
Indeed $\sigma^\epsilon=h^{-2/25}C_2^{-1/10}$, so the exact
logarithmic exponent before rounding upward is $32/5<7$.
Including \Cref{tg:good-nodes}, or using it alone when $T$ is empty,
every ratio in \Cref{tg:dual-ratio} is at most
\begin{equation}\label{tg:final-ratio}
 2h^{-4/5}+O\bigl(h^{-23/25}(1+\log h)^7\bigr)
 =o(h^{-31/40}).
\end{equation}
The exponent margins are $1/40$ and $29/200$, respectively.
Increase $h_0$ so the bound is at most $h^{-31/40}/2$.
Then $p_*\le h^{-31/40}/2$, and the defining property of the
infimum supplies a feasible positive definite $D$ with marked
averages less than $h^{-31/40}$.  No attainment assumption is used.
\end{proof}

For a fixed connected ambient graph $H$, define
\begin{equation}\label{tg:compact-domain}
 \mathcal K_H=\{X\succeq L_H:
       \ip{X}{Z}\le3\ip{L_H}{Z}\quad(Z\in\mathcal N)\}.
\end{equation}
This set is nonempty, convex, and compact.  It contains $L_H$;
closedness is immediate, and the test $\Id_{\mathsf E}$ bounds
$\tr X$ by $3\tr L_H$.  For $X\succ0$ write
$R_X(e)=b_e^TX^{-1}b_e$.  Two elementary identities used below are
\begin{equation}\label{tg:resistance-order}
 X\succeq D\succ0\ \Longrightarrow\ R_X(e)\le R_D(e),
 \qquad
 X\succeq t b_eb_e^T\ \Longleftrightarrow\ tR_X(e)\le1
 \quad(t>0).
\end{equation}
The first follows by inverse order, and the second by congruence
with $X^{-1/2}$ and the single nonzero eigenvalue of a rank-one matrix.

\begin{proposition}[Simultaneous response]\label{prop:tg-response}
Let $H$ be connected and let $G\subseteq H$ be a spanning subgraph.
Let $r$ be a sufficiently large integer, put $h=\lfloor r/2\rfloor$,
and suppose $G$ and the subgraphs induced by all parts of a nested
sequence of partitions are $h$-edge-connected, with the singleton
condition vacuous.  If $0<s\le r^{1/4}$, there is $Y\in\mathcal K_H$
such that, at every partition $\mathcal P$ in the sequence, all but
at most $r^{-1/3}|E_G(\mathcal P)|$ crossing edges satisfy
\begin{equation}\label{tg:response-edge}
 Y\succeq6s\,b_eb_e^T.
\end{equation}
Here $E_G(\mathcal P)$ denotes the edges whose endpoints lie in
different parts of $\mathcal P$.
\end{proposition}

\begin{proof}
Use the dyadic assembly of \cite[Lemma 4.2]{ThinTreesCompanion}, with
the hierarchy estimate and threshold just proved; we give the
counting argument to specify both constants.  In each occupied bin
$2^a\le |E_G(\mathcal P)|<2^{a+1}$ choose the finest partition.
List these representatives from fine to coarse as
$\mathcal P_0,\ldots,\mathcal P_{\ell-1}$, and append
$\mathcal P_\ell=\{V\}$.  Put
$B_i=E_G(\mathcal P_i)$ and $M_i=|B_i|$, with $M_\ell=0$.
Distinct parts of these representative partitions, together with $V$ and singleton leaves,
form a locally $h$-connected hierarchy.  Mark precisely the nodes
with $|O(A)|\ge|P(A)|/h$, and take $D$ from
\Cref{thm:tg-hierarchy}.  Then
\[
 Y=L_H+D\in\mathcal K_H,
\]
because $Y\succeq L_H$ and, for every $Z\in\mathcal N$,
$\ip{Y}{Z}\le\ip{L_H+L_G}{Z}\le2\ip{L_H}{Z}$.

At transition $j$ to $j+1$, the parent of each changing part is
the part of $\mathcal P_{j+1}$ containing it.  Every disappearing
edge in $B_j\setminus B_{j+1}$ belongs to parent boundaries of
changing parts, whose total full-boundary size is at most $2M_j$.
Unmarked parts therefore cover at most $2M_j/h$ such edges.
For marked parts, \Cref{tg:resistance-order} and Markov's
inequality give
\[
 |\{e\in O(A):R_Y(e)>1/(6s)\}|
 \le6s\,h^{-31/40}|O(A)|.
\]
Thus at most $2(h^{-1}+6s h^{-31/40})M_j$ disappearing edges
fail \Cref{tg:response-edge}.  Double counting only increases
this upper bound.

There is at most one representative per dyadic bin, and its bin
index strictly decreases.  Hence
$\sum_{j=i}^{\ell-1}M_j<4M_i$.  Each edge of $B_i$ disappears
at exactly one subsequent transition, so its bad-edge count is
at most $8(h^{-1}+6s h^{-31/40})M_i$.
For a skipped partition with positive crossing count $M$, its
representative is finer and has $M_i<2M$.  Consequently its
bad-edge count is at most
\begin{equation}\label{tg:response-fraction}
 16(h^{-1}+6s h^{-31/40})M
 =O(r^{-1}+r^{-21/40})M\le r^{-1/3}M
\end{equation}
for a sufficiently large absolute cutoff.  The final inequality
uses $21/40-1/3=23/120>0$.  Partitions with no crossing edges
need no estimate.  One matrix $Y$ therefore works for the whole
nested sequence.
\end{proof}
